\subsection{Case Study}

Table~\ref{tab:case-study} shows a representative failure where the agent directly queries a
complex familial relationship, leading to a tool error. \shortname identifies the root cause
as improper decomposition and induces a rule that enforces stepwise resolution of intermediate
entities. The resulting rule is abstract and reusable, and its symbolic form enables reliable
retrieval for similar queries. This example illustrates how failure traces are converted into
compact rules that correct systematic reasoning errors beyond the original instance.

\begin{table}[t]
\centering
\small
\caption{Case study illustrating failure-driven rule induction and symbolic representation.}
\label{tab:case-study}
\begin{tabular}{p{0.18\linewidth} p{0.76\linewidth}}
\toprule
\textbf{Query} &
How many letters (excluding the first and last) are there in the first name of
Viacheslav~I of Kiev's maternal grandfather? \\
\midrule
\textbf{Failure} &
The agent directly queries \textit{Viacheslav~I of Kiev's maternal grandfather}.
The tool returns an error: \textit{``no data found''}. \\
\midrule
\textbf{Proposed Rule} &
If the user query involves identifying a specific familial relationship (e.g., maternal
grandfather), then decompose the task by first resolving intermediate relationships
(e.g., mother, father) sequentially. \\
\midrule
\textbf{Symbolic Form} &
\texttt{if} \;
\texttt{domain=FAMILIAL\_RELATIONSHIP} \;\textbf{or}\;
\texttt{tool\_category=GENEALOGY\_QUERY}, \\
& \texttt{then action=[DECOMPOSE\_QUERY, RESOLVE\_INTERMEDIATE\_ENTITY, SEQUENCE\_SUBTASKS]} \\
& \texttt{strength=MANDATORY} \\
\bottomrule
\end{tabular}
\end{table}
